% Compile with XeLaTeX twice. One page for each slide of dit_results_update.tex.
\documentclass[a4paper,11pt]{article}
\usepackage{fontspec}
\usepackage[russian]{babel}
\setmainfont{Arial}
\usepackage[margin=21mm]{geometry}
\usepackage{amsmath,enumitem,hyperref}
\setlist{nosep,leftmargin=*}
\setlength{\parindent}{0pt}
\setlength{\parskip}{6pt}
\setlength{\emergencystretch}{3em}
\linespread{1.08}
\newcommand{\slide}[2]{\clearpage\section*{Слайд #1. #2}}
\newcommand{\question}[1]{\par\textbf{#1}\par}
\hypersetup{pdftitle={Понятное объяснение экспериментов по сжатию DiT},pdfauthor={},
  pdfsubject={Дано, гипотеза, действие, проверка и результат для каждого опыта},
  pdfcreator={XeLaTeX},hidelinks}
\begin{document}

\slide{1}{Как сжать модель и не испортить изображения?}
\question{Что дано?}
DiT-XL/2 --- уже обученная нейросеть. Она получает номер класса ImageNet и
создаёт изображение этого класса размером $256\times256$. Это не обучение
генератора с нуля: готовая большая модель используется как исходная точка.
Весами называются числа, определяющие вычисления сети. Их сотни миллионов,
поэтому даже хранение параметров занимает заметную память.

\question{Какая задача решается?}
Нужно хранить веса компактнее, но не потерять способность создавать хорошие
изображения. Если исходный параметр представлен 32-битным числом, а сжатое
представление требует около четырёх бит на параметр, возникает почти
восьмикратная экономия места. Это расчёт хранения, не доказательство такого
же ускорения вычислений. Ускорение зависит от реализации операций сжатой сети.

Для начала выбран VQDM --- готовый метод векторной квантизации диффузионных
моделей. Сначала он перенесён на DiT, затем проверено качество, после чего
начат поиск способов уменьшить потери. Главная линия презентации:
от обнаружения чувствительных мест к выбору обучающей цели для кодбуков.

\question{Что не является целью этих опытов?}
Не требуется доказать преимущество RVQ над VQDM. RVQ относится к ранней
истории проекта и не входит в эту презентацию. Здесь рассматриваются DiT,
перенос VQDM, его уточнение и отдельное воспроизведение VQ4DiT/DPQ.
Незавершённые измерения не выдаются за подтверждённые улучшения.

\slide{2}{Изображение получается не за один проход}
\question{Почему сеть вызывается многократно?}
Диффузионная генерация постепенно превращает случайный шум в осмысленное
изображение. В DiT этот процесс идёт не в пикселях, а в латенте: небольшом
массиве чисел, который кодирует изображение. В начале создаётся случайный
латент. Сеть получает его, время и класс и предсказывает шум, который нужно
учесть при переходе к следующему состоянию. Отдельное правило DDPM вычисляет
следующий латент. Этот новый латент снова проходит через всю ту же DiT.

В данном протоколе выполняется 50 переходов. На каждом используется одна
и та же сеть с одними и теми же весами, но разными входами и временем.
Предсказывается также дисперсия перехода: она определяет величину нового
случайного шума, добавляемого правилом DDPM. На последнем переходе эта
случайная добавка отключается. После завершения VAE-декодер превращает
итоговый латент в картинку.

\question{Чем это отличается от DDIM?}
Используемый DDPM стохастичен: случайность присутствует и в начале, и между
переходами. DDIM при нулевом параметре eta детерминирован после выбора
начального латента. Замена генератора может изменить и качество, и характер
накопления ошибки. Поэтому сравниваемые модели должны использовать один протокол.

\question{Что дало сохранение промежуточных состояний?}
Отдельный завершённый прогон сохранил состояния после 0, 10, 20, 25, 30, 40
и 50 переходов. Он позволяет увидеть постепенное появление структуры.
Это наглядная проверка работы процесса, не статистический вывод по четырём картинкам.

\slide{3}{Что находится внутри одного прохода DiT}
\question{Что означает DiT-XL/2?}
DiT --- диффузионный трансформер. Обозначение XL характеризует большую
архитектуру, а «/2» означает размер участка входного латента. Латент имеет
форму $32\times32$ и четыре канала. Участок $2\times2$ содержит
$2\cdot2\cdot4=16$ чисел. Всего участков $16\cdot16=256$.
Каждый переводится обучаемой проекцией в 1152 признака. Далее 256 векторов
совместно проходят через 28 блоков.

Attention позволяет участку учитывать информацию из других частей картинки.
Для этого используются матрица QKV и выходная матрица attention.
MLP преобразует признаки каждого участка двумя матрицами: сначала расширяет
1152 признака до 4608, затем сжимает обратно. AdaLN получает информацию
о времени и классе и выдаёт сдвиги, масштабы и управляющие множители.
В каждом блоке собственные параметры, включая собственную adaLN.

\question{Что именно сжато?}
В каждом из 28 блоков сжаты пять матриц: две attention, две MLP и одна adaLN.
Это 140 матриц, содержащих 668.86 млн весов. Веса между разными блоками
не разделяются. Но на следующем диффузионном переходе снова вызываются
все те же 28 блоков.

Выходная проекция предсказывает шум и дисперсию перехода. Она содержит
около 2.69 млн параметров и не квантуется. Также остаются несжатыми
входная проекция, представления времени/класса, смещения и VAE.
Сжатие весов не означает сжатия активаций: промежуточные данные остаются
числами с плавающей точкой.

\slide{4}{VQDM: вместо весов хранить номера в таблицах}
\question{Что такое кодбук и зачем нужен индекс?}
Кодбук --- таблица векторов. Для каждой матрицы VQDM здесь использует
четыре собственные таблицы по 256 строк. Строка содержит восемь чисел.
Восемь соседних весов приближаются суммой четырёх выбранных строк,
умноженной на масштаб соответствующей строки матрицы.
Индекс --- целый номер выбранной строки кодбука.

Вместо восьми отдельных 32-битных чисел можно хранить четыре номера от 0
до 255. Каждый номер занимает восемь бит. Получается 32 бита на группу
из восьми весов, то есть четыре бита на вес. Таблицы тоже занимают память,
но их строки многократно используются для разных групп весов. С учётом
кодбуков, масштабов и несжатых частей вся DiT занимает 4.1501 бита
на параметр в логическом упакованном представлении.

\question{Сколько всего таблиц?}
У каждой из 140 матриц свои четыре кодбука, поэтому всего их 560.
Они не общие между матрицами. В базовом варианте таблицы одной матрицы
одинаковы на всех временах генерации. Отдельные временные таблицы
появляются только в специальном опыте.

\question{Чем отличаются два способа обновления?}
Обновить кодбук --- изменить числа в одной или нескольких строках таблицы.
Это сразу влияет на все группы весов, которые используют эти строки.
Обновить индекс --- заставить конкретную группу выбрать другую строку.
Первое допускает обычный градиент; второе дискретно и требует специального поиска
или непрерывного приближения выбора.

\slide{5}{Как подобрать таблицы и проверить результат}
\question{Откуда берутся калибровочные данные?}
Готовая исходная DiT генерирует 256 траекторий по 50 переходов.
При проходе через каждый слой собирается статистика его входов.
Активациями называются значения, которые сеть вычисляет при обработке
конкретного входа. Веса определяют вычисление, активации являются данными
этого вычисления. В базовом VQDM статистика всех времён объединяется
отдельно для каждой матрицы.

\question{Что оптимизирует исходный перенос?}
При одинаковом входе нужно сблизить выход исходного и сжатого слоя.
Для линейного слоя с весами $W$ и входами $X$ ошибка имеет вид
$\|X(W-\widehat W)^\top\|_F^2/N$. Её можно вычислять через статистику
$H=X^\top X/N$. Она учитывает величины и взаимные связи входных признаков,
а не только расстояние между исходными и сжатыми весами.

Но существуют три разных проверки. Ошибка слоя сравнивает один локальный
выход. Ошибка всей DiT сравнивает предсказания двух полных сетей на одном
латенте. Ошибка генерации сравнивает финальные латенты после двух собственных
цепочек из 50 переходов. Во втором случае обе сети получают одинаковый
вход; в третьем их входы постепенно расходятся из-за предыдущих ошибок.

\question{Почему этого различия нельзя пропускать?}
Меньшая ошибка слоя не гарантирует меньшую ошибку всей сети. Более точное
предсказание сети не гарантирует меньший FID картинок. Каждый переход между
этими утверждениями требует отдельной проверки, а не только новой интерпретации
одного и того же числа.

\slide{6}{Первый результат: сжатие портит качество}
\question{Что дано и что сделано?}
Исходная DiT и её сжатая версия VQDM генерируют по 10\,000 изображений.
Условия общие: ImageNet256, DDPM50, усиление класса CFG 1.5, один VAE
и один официальный ADM-оценщик. CFG объединяет ответы сети с заданным
классом и без класса, усиливая влияние условия на генерацию.
В таблице показана 32-битная база хранения и логическая битность сжатой DiT.

\question{Что получилось?}
FID исходной DiT равен 6.192, VQDM --- 15.726. Меньше лучше.
FID оценивает различие статистик признаков реальных и сгенерированных
изображений. Он не равен доле испорченных картинок, не имеет смысла
«процент качества» и не позволяет сказать, что модель стала в определённое
число раз хуже только по отношению двух значений.

Расчётный объём весов уменьшен примерно в 7.7 раза, но качество заметно
пострадало. Следовательно, реализация сжатия работоспособна, однако
качество после прямого переноса требует улучшения. Нельзя считать,
что успешное сохранение и загрузка кодбуков уже решило задачу.

\question{Почему следующие опыты начинаются с диагностик?}
Нужно понять, какие ошибки влияют на результат. Генератор вызывается много
раз, разные слои отвечают за разные вычисления, а маленькое локальное
отклонение может усилиться позже. Диагностика временными и структурными
подстановками позволяет выбрать осмысленные гипотезы вместо случайной
смены всех настроек одновременно.

\slide{7}{Опыт 1. Когда важнее точная модель?}
\question{Что дано и какая гипотеза?}
Есть исходная DiT и готовая VQDM. Если ошибка определённой части генерации
особенно вредна, временное возвращение исходной сети на этом участке
должно сильнее приблизить итог к исходной генерации.

\question{Что сделано?}
Проверены три окна по десять переходов: начало, середина и конец.
На выбранном окне используется исходная сеть, на остальных сорока --- VQDM.
Каждый вариант продолжает собственный латент: состояние не подменяется
эталонным. Для 1000 примеров согласованы класс, начальный шум и случайные
шумы всех переходов. Поэтому расхождения не объясняются выбором другой
случайной картинки.

\question{Что измерялось?}
NMSE --- нормированная среднеквадратичная ошибка финального латента.
Для каждого примера сумма квадратов разницы делится на сумму квадратов
латента исходной модели, затем результаты усредняются. Буква N означает
normalized. Это не ошибка одного слоя и не FID.

\question{Что получилось?}
Восстановление начала уменьшило ошибку на 3.83\%, середины --- на 37.03\%,
конца --- на 1.46\%. В обратной проверке сжатая сеть используется только
в одном окне, а на остальных сорока переходах --- исходная. Середина снова
оставляет наибольшую конечную ошибку. Из трёх проверенных окон она оказалась
наиболее чувствительной. Эффекты окон нелинейны и не складываются в доли
общей ошибки. Это диагностическое смешение двух сетей, не новый способ
сохранить маленькую модель.

\slide{8}{Опыт 2. Какие части модели особенно важны?}
\question{Что изменилось относительно предыдущего опыта?}
Теперь возвращается не целая исходная сеть на части времени, а один тип
слоёв во всех блоках на всех временах. Проверяются attention, MLP и adaLN.
Остальные матрицы остаются сжатыми. Используются те же 1000 пар генераций.

Возврат attention уменьшает конечную ошибку на 11.68\%, MLP --- на 16.67\%,
adaLN --- на 30.80\%. Через adaLN в блок поступают время и класс, поэтому
её ошибка способна менять управление сразу несколькими вычислениями.
Однако adaLN составляет треть исследуемых весов. Восстановление такой
части оригинальной модели дорого и не сохраняет прежнюю битность.
Размеры семейств различны, поэтому это не сравнение чувствительности
на одинаковый восстановленный вес или бит памяти.

\question{При чём здесь шум и дисперсия?}
Дополнительно проверена подстановка выходов исходной сети на каждом вызове:
отдельно предсказания шума и отдельно дисперсии перехода. Возврат исходного
шума почти устраняет расхождение, возврат одной дисперсии влияет гораздо
слабее. Это не замена картинки после пятидесяти шагов и не возврат
новой выходной головы. Подставляется ответ сети в процессе переходов.

Возникают две проверяемые идеи. Первая --- уделять больше внимания
чувствительным временам и управляющим слоям при локальной подготовке.
Вторая --- учить сжатую сеть точнее повторять предсказание шума всей
исходной сети. Большой эффект диагностической подстановки ещё не означает,
что небольшое обучение соответствующих кодбуков даст такой же выигрыш.

\slide{9}{Опыты 3--5. Помогут ли локальные приоритеты?}
\question{Что дано и что проверялось?}
Старт --- готовый VQDM с прежней битностью. Сначала кодбуки всех матриц
дополнительно обучаются по локальному выходу слоя. Сравниваются равные веса
времён и повышенная важность средней части в два или четыре раза.
Затем в трёх выбранных блоках подгоняется выход целого блока, а не только
отдельных линейных матриц. Ещё одна группа вариантов усиливает отдельные
управляющие выходы adaLN, в частности сдвиги и множители MLP.

\question{Что получилось?}
На 256 парах генераций равномерное локальное обучение снижает конечную
ошибку на 0.65\%. Усиление середины в два раза даёт 0.96\%, в четыре ---
0.40\%. Подгонка выходов трёх блоков даёт 0.30--0.43\%, приоритет управляющих
частей adaLN --- около 0.16\%. Это небольшие эффекты, не сравнимые с
диагностическим возвратом исходной сети на среднем участке.

Два других послойных варианта не приняли ни одного обновления после
проверки округлённого FP16-представления. Дополнительные beam8-проходы
для индексов в трёх блоках дали лишь 0.09--0.70\% относительно исходного
VQDM. Сильного устойчивого усиления не получилось.

\question{Почему нельзя назвать вариант 0.96\% победителем?}
Его прямое преимущество над равномерным вариантом составляет только 0.31\%,
с интервалом от минус 0.13\% до плюс 0.87\%. Он включает ноль.
Интервал рассчитан по примерам, не по нескольким независимым обучениям.
Результат говорит о слабом эффекте этих настроек, а не о невозможности
временной калибровки вообще.

\slide{10}{Опыт 6. Свои кодбуки каждому времени?}
\question{В чём гипотеза?}
На разных этапах генерации сеть видит разные латенты. Возможно, общий
кодбук вынужден идти на компромисс. Тогда отдельные таблицы для нескольких
временных окон смогут точнее приближать веса с учётом нужных входов.
Индексы и масштабы остаются общими; для каждого окна меняются значения
в его таблицах. Это не пять полных моделей.

\question{Как сравнение сделано честнее?}
В серии с тремя окнами каждый набор обучается 200 обновлений, общий контроль
получает 600. Окна содержат первые 30, следующие 10 и последние 10 переходов.
В серии с пятью равными окнами по 10 переходов каждый набор получает 200,
контроль --- 1000.
Для обоих вариантов внутри каждой серии оцениваются 10\,000 изображений.
Сравнивается не разное количество обучения одного и того же общего набора,
а распределение сопоставимого бюджета между таблицами.

\question{Что получилось?}
Три набора дали FID 15.663 против 15.536 у общего. Пять наборов --- 15.695
против 15.644. Для пяти наборов добавляется около 9.18 МБ хранения.
Локальная ошибка собственного окна снижается на 1.27--3.96\%, но FID
не становится лучше общего контроля. Малые различия не объявляются
статистически значимыми без повторных обучений.

Отдельная проверка трёх наборов показала относительное различие значений
около 1.3--1.8\%; направления маленьких поправок разные. Близость таблиц
сама по себе не доказывает их взаимозаменяемость, но отдельные измерения
генерации пока не оправдали дополнительную память.

\slide{11}{Другая идея: повторять ответ всей исходной сети}
\question{Чем это отличается от локального обучения?}
В локальном подходе слой учится повторять выход исходного слоя на сохранённых
входах. В новой постановке один и тот же зашумлённый латент, время и класс
подаются двум полным сетям. Исходная сеть выдаёт ответ-образец и называется
учителем. Сжатая сеть тоже проходит все 28 блоков. Сравнивается уже их
итоговое предсказание шума.

\question{Как обновляются кодбуки?}
Средний квадрат разницы ответов задаёт обучающую ошибку. Обратное
распространение показывает, как изменение каждого числа в кодбуках
влияет на эту ошибку через всю сжатую сеть. Оптимизатор немного меняет
значения кодовых векторов. Индексы не меняются, исходная DiT не обучается,
другие параметры сохраняются. Каждое изменение кодового вектора влияет
на все группы весов, которые используют соответствующий индекс.

Ожидаемая польза в том, что ошибки разных слоёв могут компенсироваться.
Их не обязательно сделать по отдельности минимальными, если совместный
результат всей сети стал ближе к исходному. Это дистилляция --- обучение
по ответам уже умеющей решать задачу модели.

\question{Обучение разворачивает все 50 переходов?}
Нет. Для одного обучающего примера выполняется один вызов исходной DiT
и один вызов сжатой. Градиент проходит через 28 блоков сжатой сети,
но не через цепочку полной генерации. Пользу для генерации затем
нужно проверять отдельно. Это дополнительное обучение кодбуков,
поэтому называть такую процедуру training-free некорректно.

\slide{12}{Опыт 7. Достаточно ли подстроить три блока?}
\question{Что дано и что изменено?}
Старт --- готовый VQDM. Для быстрого пилота обучаются кодбуки только
4-го, 15-го и 26-го блоков, то есть 15 матриц. Другие кодбуки, индексы,
масштабы и смещения заморожены. За 200 обновлений ответы сравниваются
после всей сети, несмотря на обучение лишь её части.
Входы взяты из пяти моментов генерации исходной DiT.

Проверены три цели: обычное предсказание шума; ответ после усиления класса
CFG; тот же ответ, но средние времена имеют вес в четыре раза больше.
Это проверка не только чувствительного места, но и уровня обучающего сигнала.

\question{Что получилось?}
При полной генерации 256 пар примеров исходная ошибка VQDM равна 0.136118.
После обучения по обычному шуму она стала 0.130247, то есть ниже на 4.31\%.
По ответу CFG --- 0.129780, снижение 4.66\%. С усилением середины ---
0.128446, снижение 5.64\%. Для всех трёх сравнений с исходным VQDM
95\%-е интервалы улучшения не содержат ноль.

\question{Какой вывод допустим?}
Обучение по выходу всей DiT уменьшило не только обучающую ошибку,
но и расхождение после 50 переходов. Это положительный сигнал.
Однако результат не равен улучшению FID. Кроме того, дополнительная
польза усиленной середины над обычным CFG отдельно не установлена:
интервал прямого сравнения включает ноль. Интервалы отражают разброс
по примерам при фиксированном обучении, а не устойчивость при смене его сида.

\slide{13}{Опыт 8. Какой ответ использовать для обучения?}
\question{Почему это отдельная серия?}
Здесь старт не готовый VQDM. Кодбуки только инициализированы по весам
с помощью остаточного k-means. Это более слабая и другая начальная точка.
Обучаются все 140 матриц на 4000 латентах реальных изображений ImageNet.
Ещё 1000 изображений отложены; контрольная ошибка считается на одинаковых
96 сочетаниях изображения и времени.

\question{Какие цели сравниваются?}
Первая просит сжатую сеть повторять выходы всех 28 блоков учителя.
Вторая сравнивает итоговое предсказание шума всей DiT.
Третья сравнивает предсказание с известным случайным шумом, который
был добавлен к чистому латенту. Во всех трёх вариантах дополнительно
используется небольшая ошибка дисперсии относительно исходной сети.
Индексы в этих трёх ветвях фиксированы.

\question{Что получилось?}
Общая ошибка относительно ответа исходной сети до обучения равна 0.003274.
На шаге 200 подгонка выходов блоков даёт 0.002026, ответа всей сети ---
0.001734, известного шума --- 0.003015. Лучше сработало повторение ответа
всей DiT: снижение ошибки на 47.04\%. На шаге 1000 у него уже 0.002185.
Продолжение обучения ухудшило контрольный показатель.

Это мотивация выбирать момент остановки по отложенным данным, а не считать,
что больше шагов всегда лучше. Процент 47.04 относится только к этой
инициализации и этой ошибке. Он не является выигрышем над готовым VQDM
и не заменяет оценку качества изображений.

\slide{14}{Опыт 9. Поможет ли ещё и менять индексы?}
\question{Почему индексы вообще хотелось менять?}
Если таблицы сдвинулись во время обучения, прежний номер строки может
перестать быть лучшим для конкретной группы весов. Поэтому возникает
идея чередовать обычные обновления кодбуков и перевыбор некоторых индексов.
Выбор номера дискретен: градиент не может непосредственно изменить целое
число от 0 до 255. В этом опыте используется явная проверка кандидатов.

\question{Что сделано?}
Взяты та же свежая инициализация и три цели предыдущего слайда.
Для каждой цели запущен второй вариант: раз в 25 обновлений проверяются
редкие замены индексов. В событии обрабатываются четыре слоя,
выбираются группы с большим сигналом градиента, допускается до 128
изменений на слой. Замена принимается только после проверки снижения
реальной учебной ошибки; при необходимости выполняется откат кандидата.

\question{Что получилось?}
На отложенных входах после 200 обновлений все три варианта хуже контроля
с фиксированными индексами. Для выходов блоков ошибка 0.002481 вместо
0.002026, для ответа всей сети 0.002373 вместо 0.001734, для известного
шума 0.004133 вместо 0.003015. Поздние точки не исправили общий вывод.

Следовательно, уменьшать ошибку на выбранной учебной проверке оказалось
недостаточно для улучшения отложенного результата. Но это очень разреженная
схема: принято лишь около 4--11 тысяч изменений из 334 млн назначений,
возможны повторные изменения одного места. Она не опровергает другие
способы глобальной оптимизации дискретных индексов.

\slide{15}{Опыт 10. Все кодбуки готового VQDM}
\question{Что дано?}
После пилотов проверяется основная практическая постановка: улучшить
готовую сжатую VQDM, а не только слабую стартовую инициализацию.
Все 140 матриц уже подготовлены с учётом активаций. Их структура остаётся
той же: четыре кодбука по 256 векторов длины восемь на матрицу.

\question{Какие входы используются?}
4000 реальных изображений ImageNet переводятся VAE в чистые латенты.
Для обучающего примера выбирается случайное время из исходных 1000
уровней, добавляется соответствующий шум. Это прямое зашумление реального
латента, а не запуск 50 переходов генерации для каждого обновления.
Полученный вход, время и класс одинаковы для исходной и сжатой DiT.

\question{Что обучается?}
Только значения кодбуков всех матриц. Индексы, масштабы строк, смещения
и остальные параметры фиксированы. Минимизируется ошибка предсказания шума
относительно исходной сети плюс ошибка сырого выхода дисперсии с весом 0.1.
Выполняются 500 обновлений Adam с шагом $10^{-5}$ и эффективным батчем 16.
При обучении CFG не применяется, класс иногда отключается.

Для контроля отложены 1000 изображений. На фиксированных 96 входах,
по 16 на каждом из шести времён, проверяется ошибка ответа всей сети.
Эти входы не используются для обучения. Битность остаётся 4.1501:
таблицы меняют содержимое, но их количество и размеры не растут.
Это дополнительная оптимизация кодбуков через сеть, а не новый исходный
обучающий прогон всех полноточных весов DiT.

\slide{16}{Результат: сеть точнее повторяет исходную}
\question{Как читать график?}
По горизонтали отложено число обновлений кодбуков: от нуля до пятисот.
Это не время диффузии и не номер перехода генерации. По вертикали ---
средний квадрат разницы между предсказаниями шума двух полных сетей
на одинаковых контрольных входах. Меньше означает точнее повторяется
ответ исходной DiT. Контрольные входы не меняются между точками кривой.

До обучения ошибка равна 0.001162851. На шаге 50 --- 0.001087955,
на 100 --- 0.001061323, на 200 --- 0.001039287,
на 500 --- 0.001007803. Суммарное относительное снижение равно 13.33\%.
В этом опыте контрольный показатель продолжал улучшаться до последнего
проверенного шага. График не начинается от нуля по вертикали, поэтому
важно читать подписи чисел и процент, а не оценивать улучшение по высоте линии.

\question{Почему здесь не 47\%, как раньше?}
Потому что раньше была свежая инициализация по весам, а здесь готовая
activation-aware VQDM с меньшей стартовой ошибкой. Сравниваются разные
исходные состояния, шаг обучения тоже различается. Проценты нельзя
приписывать одной и той же модели или складывать между собой.

\question{Что этот результат доказывает и чего не доказывает?}
Сжатая сеть точнее повторяет учителя на отложенных входах, без роста
битности и без смены индексов. Это завершённый положительный результат.
Но он ещё не доказывает снижение FID, устойчивость к независимому сиду
обучения или преимущество над методами из других статей.

\slide{17}{Опыты 11--14. Дополнения к подстройке}
\question{Что является контролем?}
Простое обучение кодбуков готового VQDM по ответу учителя, как на предыдущих
двух слайдах. Каждый дополнительный вариант начинает с того же исходного
VQDM и имеет ту же битность. Для оценки получены 5000 общих изображений,
а сравнения FID выполнены внутри этого единого протокола.

Первая добавка обучает масштабы строк матриц вместе с кодбуками.
FID равен 11.050 против 11.045 у контроля: убедительного улучшения нет.
Вторая использует смесь входов: половина --- обычные зашумлённые реальные
латенты, половина --- состояния заранее сохранённых генераций начальной
сжатой сети. FID 11.146, то есть не лучше контроля. Это не обучение
на непрерывно обновляемых собственных траекториях; дополнительно меняются
особенности выборки времени и классов.

Третья заменяет ответ учителя на известный шум и полноценную DDPM-цель
дисперсии. Лучший по контрольной ошибке шаг 100 даёт FID 13.390,
финальный 500 --- 14.350. Выбор сделан до оценки FID, не по картинкам.
Это адаптация цели, а не воспроизведение всего DPQ.

Четвёртая позволяет плавно выбирать между текущим и альтернативным
индексом, а при оценке оставляет один номер. Изменено 3.92\% назначений,
ошибка ответа выросла в 5.70 раза, FID достиг 24.376. Эта схема не помогла
и не является репродукцией VQ4DiT. Отсутствие выигрыша относится к FID:
другие метрики у некоторых вариантов немного лучше, но общий основной
кандидат остаётся простым обучением кодбуков по учителю.

\slide{18}{Воспроизвести VQ4DiT и DPQ}
\question{Зачем нужны отдельные репродукции?}
Можно неудачно перенести одну идею статьи, оставив другие компоненты
чужой системы. Это не проверка самой статьи. Поэтому параллельно
подготовлены самостоятельные воспроизведения VQ4DiT и DPQ на DiT-XL/2.
Их результаты не смешиваются с FID адаптаций предыдущего слайда.

В VQ4DiT каждая матрица получает свой кодбук, а назначение вектора
уточняется через обучаемый выбор кандидатов. Запуск закреплён на авторском
коде с отдельным протоколом расхождений между выпуском кода и статьёй.
Использованы 500 обновлений, кодбуки и смещения обучаются, мягкий выбор
постепенно превращается в дискретный. В сохранённом срезе обучение
завершилось, но окончательные метрики изображений ещё отсутствуют.

Для DPQ подготовлены опубликованный вариант и отдельное восстановление
недостающих частей описания статьи. Они обучены по одной эпохе на полном
ImageNet. Обучение обоих завершено, генерация для оценки ещё не завершена
в сохранённом состоянии. Одно завершённое обучение не считается
подтверждением качества из статьи.

\question{Почему нельзя сразу поставить все FID в один рейтинг?}
У DPQ используется DDPM250, у VQ4DiT --- DDPM50 и другой VAE.
Различаются также битность и способ учёта кодбуков. Нужны собственные
контроли исходной DiT в тех же условиях и явное сопоставление протокола.
Пока репродукции являются завершёнными обучениями с незавершённой оценкой,
а не положительными или отрицательными результатами по FID.

\slide{19}{Что получилось установить}
\question{Как связаны все группы опытов?}
Первый перенос VQDM дал работающую сжатую DiT, но ухудшил FID.
Затем диагностические замены показали, что средний участок генерации
и adaLN заметно влияют на конечное расхождение. Эти наблюдения мотивировали
локальные приоритеты и временные таблицы. Однако такие дополнения дали
небольшие эффекты и не оправдали усложнение по измеренному качеству.

После этого обучающий сигнал перенесён с отдельных слоёв на ответ всей
исходной сети. Пилот на трёх блоках уменьшил конечную ошибку генерации
на 4.3--5.6\%. Отдельное обучение кодбуков всех матриц готового VQDM
уменьшило ошибку предсказания шума на отложенных входах на 13.33\%.
Свежая инициализация дополнительно показала важность цели и остановки
обучения: хорошая промежуточная точка может оказаться лучше поздней.

\question{Что здесь главное?}
Чувствительность части системы не даёт автоматически готового улучшения.
Лучший текущий сигнал связан с тем, какое поведение просится повторять:
полный ответ исходной DiT оказался полезнее набора локальных поправок.
При этом положительные проценты описывают разные проверки и разные
начальные модели; их нельзя складывать или выдавать за уменьшение FID.

Проверенные изменения индексов, масштабов, распределения входов и DDPM-цели
не улучшили FID простого контроля. Это ограничивает приоритет следующих
опытов, но не доказывает невозможности их иных реализаций.
Самый разумный следующий шаг --- подтвердить основной результат на
качестве изображений и независимых повторах, а не добавлять сразу
несколько новых механизмов.

\slide{20}{Что нужно, прежде чем заявлять улучшение}
\question{Какого результата ещё нет?}
Новый FID простого обучения равен 11.045 на 5000 изображений.
Старый FID готового VQDM равен 15.726 на 10\,000 изображений.
Число изображений влияет на оценку FID, поэтому эти значения нельзя
использовать как честную пару «до и после». Это не формальность:
сам переход между объёмами выборки способен изменить число без изменения модели.

Нужен готовый VQDM до подстройки на тех же 5000 примерах, с теми же
случайными числами, генератором, VAE, эталонными статистиками и оценщиком.
Такой контроль включён в отдельную оценку, но её финальные метрики
отсутствуют в сохранённом срезе. Генерация файлов или состояние задачи
в очереди сами по себе не заменяют завершённый отчёт метрик.

\question{Что делать после честного сравнения?}
Если выигрыш подтверждается, повторить обучение с независимыми сидами
и проверить небольшой набор шагов обучения и моментов остановки.
Контрольную точку выбирать по отложенным данным до просмотра итогового FID.
Далее оценить выбранный вариант на более крупном общем наборе.
Одновременно завершить оценки VQ4DiT и DPQ и сопоставить их с соответствующими
статьями, не игнорируя различия генераторов и VAE.

Рабочая гипотеза остаётся конкретной: компактные таблицы могут лучше
сохранять генеративное поведение, если обучаются повторять ответ всей
исходной сети. Для её подтверждения уже есть ошибки сети и генерации,
но ещё требуется согласованный результат по качеству изображений.

\end{document}
